% Generated draft: requires mathematical and source review.
\documentclass[11pt]{article}
\usepackage[margin=25mm]{geometry}
\usepackage{fontspec}
\setmainfont{Latin Modern Roman}
\usepackage{amsmath,amssymb,mathtools,unicode-math}
\setmathfont{Latin Modern Math}
\usepackage{enumitem,needspace,xurl,hyperref}
\setlength{\parindent}{0pt}
\setlength{\parskip}{5pt}
\setlength{\emergencystretch}{3em}
\newcommand{\N}{\mathbb{N}}
\newcommand{\Z}{\mathbb{Z}}
\newcommand{\Q}{\mathbb{Q}}
\newcommand{\R}{\mathbb{R}}
\newcommand{\C}{\mathbb{C}}
\newcommand{\F}{\mathbb{F}}
\newcommand{\A}{\mathbb{A}}
\newcommand{\PP}{\mathbb{P}}
\newcommand{\E}{\mathbb{E}}
\newcommand{\eps}{\varepsilon}
\newcommand{\dd}{\,\mathrm{d}}
\begin{document}
\section*{UnsolvedMath ID: 30000971}
\textbf{Source problem code:} catalogue entry
\subsection{Definitions and mathematical statement}
Let \(X\subseteq\PP^r\) be a smooth projective variety of dimension \(n\), and let \(c\geq1\) be an integer with \(n+c\leq r\). Let \(\pi:X\to\PP^{n+c}\) be a general linear projection, with center disjoint from \(X\), and let \(Z_y=\pi^{-1}(y)\) be its scheme-theoretic fiber. The field, geometric-point convention, and precise meaning of general are not specified in the extraction.

Eisenbud's fiber-regularity conjecture as supplied asks whether
\[
\operatorname{reg}(Z_y)\leq\frac nc+1
\quad\text{for every }y\in\PP^{n+c}.
\]
The regularity is an integer, so when defined this bound is equivalently \(\operatorname{reg}(Z_y)\leq\lfloor n/c\rfloor+1\). The statement is simultaneous for every fiber of a general projection; it does not merely concern a general fiber.

The same record gives a second conjectural inequality involving the obstruction sheaf \(Q_y\) attached to the fiber:
\[
\operatorname{reg}(Z_y)\leq\frac1c\operatorname{length}(Q_y).
\]
Length here requires the finite-length obstruction object intended in the report, whose definition is absent. No substitute such as fiber length, tangent-space dimension, or normal sheaf is inserted for \(Q_y\).

For clarity, one common convention for the regularity of a projective subscheme \(Z\subseteq\PP^r\) is the Castelnuovo--Mumford regularity of its ideal sheaf, the least integer \(m\) for which
\[
H^i(\PP^r,\mathcal I_Z(m-i))=0\qquad(i>0).
\]
The source does not fix this convention versus regularity of the homogeneous coordinate ring or structure sheaf, which can shift the relevant integer. It also quantifies over points outside the image, where the fiber is empty, without specifying empty-fiber conventions. These details and the definition of \(Q_y\) are necessary before either inequality is fully self-contained.
\subsection{Short English statement}
For a general projection of a smooth projective \(n\)-fold to \(\PP^{n+c}\), is every scheme fiber at most \(n/c+1\) in regularity? A second question bounds that regularity by the obstruction-sheaf length divided by \(c\).
\subsection{Sources}
\begin{itemize}
\item Local UnsolvedMath registry, permanent ID \(30000971\), source code catalogue entry, ``Eisenbud's Fiber-Regularity Conjecture.'' Attribution to Eisenbud is explicit in the original statement; no individual proposer field is populated.
\item \emph{Classical Algebraic Geometry} (2009), Oberwolfach Reports, DOI 10.4171/owr/2008/26, pp. 1425--1474: \url{https://doi.org/10.4171/owr/2008/26}.
\item Supplied reference, ``Beheshti--Eisenbud, Fibers of generic projections'': \url{https://arxiv.org/abs/0806.1928}.
\item Supplied literature triage, checked 2026-08-21: status ``open.'' It identifies the first target as Eisenbud's fiber-regularity conjecture, reports important bounds and reductions by Beheshti and Eisenbud, and calls the general uniform bound unresolved. This dated assessment derives from supplied OWR curation; time-sensitive status and cited claims require independent verification.
\end{itemize}
\end{document}
